% ============================================================================
\chapter{Interface moderne : Vite, composants, React/Vue/Svelte, Astro}
\label{ch-js-frameworks}
% ============================================================================

\begin{objectif}
\begin{itemize}
  \item D\'emarrer un projet avec Vite et comprendre ce qu'un bundler fait pour toi.
  \item Saisir le mod\`ele \og ~\'etat $\rightarrow$ rendu~ \fg{} commun aux frameworks, et les hooks React.
  \item Choisir un cadre en connaissance de cause, sans religion.
  \item G\'erer \'etat, routes, donn\'ees serveur et styles dans une vraie application.
  \item Savoir quand \emph{ne pas} mettre de framework.
\end{itemize}
\end{objectif}

\section{Vite~: le socle}

\link{Vite}{https://vitejs.dev/} (7\textsuperscript{e}~majore en 2025, cr\'e\'e par Evan~You) sert tes
modules nativement en ESM pendant le d\'eveloppement (d\'emarrage quasi instantan\'e, HMR cibl\'e) et
utilise \link{Rollup}{https://rollupjs.org/} pour construire la version de production
(minification, d\'ecoupage de code, empreintes de fichiers).

\begin{terminal}
npm create vite@latest site -- --template vanilla        # ou react, vue, svelte, preact, lit...
cd site
npm install
npm run dev              # http://localhost:5173
npm run build && npm run preview
\end{terminal}

\begin{codejs}
// vite.config.js --- les reglages qui reviennent vraiment
import { defineConfig } from "vite";
import react from "@vitejs/plugin-react";

export default defineConfig({
  plugins: [react()],
  resolve: { alias: { "@": new URL("./src", import.meta.url).pathname } },
  server: { port: 5173, host: true, proxy: { "/api": "http://localhost:8000" } },
  build: {
    sourcemap: true,
    target: "es2022",
    rollupOptions: { output: { manualChunks: (id) => (id.includes("node_modules") ? "deps" : undefined) } },
  },
  define: { __APP_VERSION__: JSON.stringify(process.env.npm_package_version) },
});
\end{codejs}

Les variables d'environnement expos\'ees au client doivent \^etre pr\'efix\'ees \verb|VITE\_|~:
\verb|import.meta.env.VITE\_API\_URL|. Tout le reste du fichier \verb|.env| reste c\^ot\'e build~:
c'est intentionnel, ne cherche pas \`a le contourner.

\section{Le mod\`ele commun \`a tous les frameworks}

\begin{enumerate}
  \item Tu d\'ecris l'\textbf{\'etat} (donn\'ees) et une \textbf{fonction de rendu}
        (\verb|\'etat -> HTML|).
  \item Tu changes l'\'etat avec les moyens du framework (\verb|setState|, \verb|\$set|, signal,
        assignation r\'eactive).
  \item Le framework \textbf{diff} l'ancien et le nouveau rendu, applique le minimum
        d'op\'erations DOM (la \og ~r\'econciliation~ \fg{}), puis compose.
\end{enumerate}

Tout le reste (hooks, signaux, stores, effets) n'est que de la broderie autour de ce mod\`ele.
C'est pour cela qu'apprendre un framework en rend les autres rapides.

\subsection{React --- le cas le plus demand\'e}

\begin{codejs}
// src/Compteur.jsx --- lecture obligatoire : un composant est une fonction pure de l'etat
import { useState, useEffect, useMemo, useCallback, useRef, useReducer } from "react";

const initial = () => JSON.parse(localStorage.getItem("courses") ?? "[]");

export function ListeCourses() {
  const [articles, setArticles] = useState(initial);          // initialise une seule fois
  const [filtre, setFiltre] = useState("tous");
  const [brouillon, setBrouillon] = useState("");
  const champ = useRef(null);                                  // reference a un noeud DOM, sans re-rendu

  useEffect(() => {                                            // effet = synchronisation avec l'exterieur
    localStorage.setItem("courses", JSON.stringify(articles));
  }, [articles]);

  useEffect(() => {                                            // ecouteur global + nettoyage
    const onKey = (e) => e.key === "/" && champ.current?.focus();
    window.addEventListener("keydown", onKey);
    return () => window.removeEventListener("keydown", onKey);
  }, []);                                                        // [] : une fois, au montage

  const restants = useMemo(() => articles.filter((a) => !a.faite), [articles]);
  const total = useMemo(() => restants.reduce((s, a) => s + a.qte, 0), [restants]);

  const ajouter = useCallback(() => {
    const nom = brouillon.trim();
    if (!nom) return;
    setArticles((liste) => [{ id: crypto.randomUUID(), nom, qte: 1, faite: false }, ...liste]);
    setBrouillon("");
  }, [brouillon]);

  const basculer = (id) =>
    setArticles((liste) => liste.map((a) => (a.id === id ? { ...a, faite: !a.faite } : a)));

  return (
    <main>
      <h1>Courses <small>({total})</small></h1>
      <form onSubmit={(e) => { e.preventDefault(); ajouter(); }}>
        <input ref={champ} value={brouillon} onChange={(e) => setBrouillon(e.target.value)}
               placeholder="Ajouter (touche /)" />
      </form>
      <ul>
        {restants.map((a) => (
          <li key={a.id}>
            <label>
              <input type="checkbox" checked={a.faite} onChange={() => basculer(a.id)} />
              {a.nom}
            </label>
          </li>
        ))}
      </ul>
    </main>
  );
}
\end{codejs}

\etape{Les quatre r\`egles des hooks, sans exception}
\begin{enumerate}
  \item Un hook s'appelle \verb|useXxx|, et uniquement au \textbf{niveau sup\'erieur} d'un
        composant ou d'un hook~: jamais dans une condition, une boucle, un rappel.
  \item Un \verb|useEffect| doit d\'eclarer ses d\'ependances \textbf{compl\`etes}
        (active \verb|eslint-plugin-react-hooks|, et le preset \og ~recommended-latest~ \fg{}).
  \item Une mise \`a jour d'\'etat doit produire une \textbf{nouvelle r\'ef\'erence}
        (\verb!set(x => ({ ...x }))!, pas \verb!set(state.prop = 1)!).
  \item Si un effet sert \`a \og ~d\'eriver~ \fg{} une valeur, ce n'est pas un effet~: c'est un calcul
        pendant le rendu (\verb|useMemo| ou simple const).
\end{enumerate}

\subsection{Vue~3 (composition API)}

\begin{codeplain}
<!-- ArticleList.vue -->
<script setup>
import { ref, computed, watchEffect } from "vue";

const articles = ref(JSON.parse(localStorage.getItem("courses") ?? "[]"));
const brouillon = ref("");
const restants = computed(() => articles.value.filter((a) => !a.faite));

function ajouter() {
  const nom = brouillon.value.trim();
  if (!nom) return;
  articles.value.unshift({ id: crypto.randomUUID(), nom, qte: 1, faite: false });
  brouillon.value = "";
}
watchEffect(() => localStorage.setItem("courses", JSON.stringify(articles.value)));
</script>

<template>
  <input v-model="brouillon" @keyup.enter="ajouter" placeholder="Article" />
  <ul><li v-for="a in restants" :key="a.id">{{ a.nom }}</li></ul>
</template>
\end{codeplain}

\subsection{Svelte~5 (runes)~: le moins de code du march\'e}

\begin{codeplain}
<script>
  let articles = $state(JSON.parse(localStorage.getItem("courses") ?? "[]"));
  let brouillon = $state("");
  const restants = $derived(articles.filter((a) => !a.faite));
  $effect(() => localStorage.setItem("courses", JSON.stringify(articles)));

  function ajouter() {
    const nom = brouillon.trim();
    if (!nom) return;
    articles = [{ id: crypto.randomUUID(), nom, qte: 1, faite: false }, ...articles];
    brouillon = "";
  }
</script>

<input bind:value={brouillon} on:keydown={(e) => e.key === "Enter" && ajouter()} />
<ul>{#each restants as a (a.id)}<li>{a.nom}</li>{/each}</ul>
\end{codeplain}

\section{Panorama honn\^ete}

\begin{tabz}{lXlX}
  \headrow \thd{Cadre} \& \thd{Id\'ee forte} \& \thd{Taille (gzipp\'ee, app. vide)} \& \thd{Doc} \\
  \midrule
  \link{React~19}{https://react.dev/} \& le plus d'emplois, \'ecosyst\`eme immense, rendu par fonction \& $\approx$45~Ko \& \url{https://react.dev/learn} \\
  \link{Vue~3.5}{https://fr.vuejs.org/} \& la meilleure mont\'ee en charge du d\'ebutant au projet moyen, templates clairs \& $\approx$35~Ko \& \url{https://fr.vuejs.org/guide/introduction.html} \\
  \link{Svelte~5 / SvelteKit}{https://svelte.dev/} \& compilateur, peu de code, performances par d\'efaut \& $\approx$7~Ko \& \url{https://learn.svelte.dev/} \\
  \link{Solid}{https://www.solidjs.com/} \& signaux fins, syntaxe JSX, rapidit\'e maximale \& $\approx$8~Ko \& \url{https://docs.solidjs.com/} \\
  \link{Qwik}{https://qwik.dev/} \& reprise d'hydratation~: ne charge que ce que l'utilisateur touche \& minuscule \& \url{https://qwik.dev/learn/} \\
  \link{Angular~19+}{https://angular.dev/} \& tout int\'egr\'e (routes, HTTP, formulaires, DI), signals, standalone \& lourd mais complet \& \url{https://angular.dev/tutorials} \\
  \link{Astro~5}{https://astro.build/} \& sites centr\'es contenu~: z\'ero JS par d\'efaut, \^ilots d'interactive \& 0~Ko \& \url{https://docs.astro.build/fr/getting-started/} \\
  \link{HTMX}{https://htmx.org/} \& l'interactive dans les attributs HTML, le serveur renvoie du HTML \& 14~Ko \& \url{https://htmx.org/docs/} \\
  \link{Lit}{https://lit.dev/} \& web components natifs et durables, multi-cadres \& $\approx$17~Ko \& \url{https://lit.dev/docs/} \\
  \link{Alpine.js}{https://alpinejs.dev/} \& du \og ~jQuery propre~ \fg{} pour les petites pages \& $\approx$15~Ko \& \url{https://alpinejs.dev/start-here} \\
  \link{jQuery}{https://jquery.com/} \& \`a conna\^itre pour maintenir l'existant, pas pour d\'emarrer \& 30~Ko \& \url{https://api.jquery.com/} \\
\end{tabz}

\subsection{Les briques transverses}

\begin{tabz}{lXl}
  \headrow \thd{Besoin} \& \thd{Choix raisonnable en 2026} \& \thd{Lien} \\
  \midrule
  Routage SPA \& TanStack Router (typ\'e) ou react-router v7 \& \url{https://tanstack.com/router} \\
  Donn\'ees serveur \& TanStack Query~: cache, retries, invalidation, \'etats \& \url{https://tanstack.com/query} \\
  \'Etat client \& Zustand (l\'eger) ou le signal natif ; Redux Toolkit si grosse \'equipe \& \url{https://zustand.docs.pmnd.rs/} \\
  Formulaires \& React Hook Form + zod (validation partag\'ee client/serveur) \& \url{https://react-hook-form.com/} \\
  Styles \& CSS natif (variables, \verb|:has|, \verb|@layer|, \verb|container| queries) ou Tailwind~4 \& \url{https://tailwindcss.com/} \\
  Composants UI \& shadcn/ui (copie le code, ne d\'epend pas d'une lib) \& \url{https://ui.shadcn.com/} \\
  Animation \& \verb|element.animate()|, Motion (framer-motion), GSAP \& \url{https://motion.dev/} \\
  Graphes/donn\'ees \& \link{D3}{https://d3js.org/}, \link{ECharts}{https://echarts.apache.org/}, \link{Chart.js}{https://www.chartjs.org/} \& \\
\end{tabz}

\section{Quand ne pas mettre de framework}

\begin{itemize}
  \item Site vitrine, blog, documentation, landing page~: \textbf{Astro} (ou du HTML + un peu de JS).
  \item Page avec trois interactions~: HTML + 80 lignes de JS natif + \link{Alpine}{https://alpinejs.dev/}.
  \item Backend qui rend du HTML (PHP, Java, Rails, Laravel)~: \link{HTMX}{https://htmx.org/}
        ou \link{Hotwire/Turbo}{https://turbo.hotwired.dev/}~: tu gardes la simplicit\'e du serveur,
        tu gagnes le confort du SPA. C'est une tendance lourde et saine.
  \item Composant autonome r\'eutilisable dans n'importe quelle page~: un
        \link{Web Component}{https://developer.mozilla.org/fr/docs/Web/API/Web\_components} natif.
\end{itemize}

\begin{note}
Un framework se choisit sur trois crit\`eres~: la taille de l'\'equipe, la part d'interactive r\'eelle
de l'interface, et la disponibilit\'e des comp\'etences sur le march\'e local. Pour un portfolio,
React ou Vue maximisent les chances d'entretien~; pour un projet solo qui doit sortir vite,
Svelte ou Astro gagnent.
\end{note}

\begin{exercice}
\begin{enumerate}
  \item Initialise trois projets du m\^eme contenu (liste de courses)~: vanilla, Vue, React.
        Compare le nombre de lignes, le poids du bundle (\verb|dist/|) et le temps de
        d\'eveloppement~: \'ecris ton avis en dix lignes dans un \verb|ANALYSE.md|.
  \item Sur le projet React~: remplace le \verb|useState| local par un store \verb|zustand|,
        et ajoute une persistance via le middleware \verb|persist|.
  \item Branche \link{TanStack~Query}{https://tanstack.com/query} sur une API r\'eelle~: g\`ere
        \'etats \og ~chargement / erreur / vide~ \fg{}, mise en cache, invalidation apr\`es mutation.
  \item Ajoute le routage~: routes \verb|/articles|, \verb|/articles/:id|, \verb|404|,
        et un chargement diff\'er\'e (\verb|React.lazy| / \verb|defineAsyncComponent|) de la page d\'etail.
  \item Migre un composant React existant vers un Web Component natif (\verb|customElements.define|)
        et utilise-le dans une page HTML sans framework.
\end{enumerate}
\end{exercice}

\begin{ressources}
\begin{itemize}
  \item Vite --- \url{https://vite.dev/guide/}
  \item React (nouvelle doc, excellente) --- \url{https://react.dev/learn}
  \item Vue --- \url{https://fr.vuejs.org/guide/extras/composition-api-faq.html}
  \item Svelte --- \url{https://svelte.dev/tutorial/} --- Astro --- \url{https://docs.astro.build/fr/core-concepts/overview/}
  \item State of JS (sondage annuel, pour choisir sans se faire d'opinion fausse) --- \url{https://2024.stateofjs.com/}
\end{itemize}
\end{ressources}

\index{Vite}\index{React}\index{Vue}\index{Svelte}\index{framework}
